\section{Architecture}
Notre implémentation des threads repose sur une structure appelée \textit{thread\_struct}. Nous allons expliciter ici la structure élémentaire définissant un thread.\\ 
Certains éléments avancés de cette structure ont été volontairement omis et seront précisés plus tard dans ce document.

\begin{itemize}
    \item \texttt{state} : état du thread.
    \item \texttt{ctx} : contexte d’exécution.
    \item \texttt{stack} : pointeur vers la pile allouée pour ce thread.
    \item \texttt{func} : fonction exécutée par le thread
    \item \texttt{funcarg} : arguments de la fonction du thread
    \item \texttt{retval} : valeur de retour de la fonction.
\end{itemize}

Comme exprimé précédemment, cette structure est incomplète. Nous allons désormais aborder les problématiques nous ayant poussé à modifier \textit{thread\_struct}.






